\documentclass[11pt]{amsart}

\usepackage[T1]{fontenc}
\usepackage{lmodern}
\usepackage{microtype}
\usepackage{amsmath,amssymb,amsthm}
\usepackage[hidelinks]{hyperref}

\newtheorem{theorem}{Theorem}
\newtheorem{lemma}[theorem]{Lemma}
\theoremstyle{definition}
\newtheorem{remark}[theorem]{Remark}

\DeclareMathOperator{\Aut}{Aut}

\title[The Liu--Li predicted degree sequence at $(8,5)$]{The Liu--Li Predicted
Degree Sequence Is Not Extremal at
\texorpdfstring{$(n,\Delta)=(8,5)$}{(n,Delta)=(8,5)}}
\date{}
\subjclass[2020]{Primary 05C50; Secondary 05C35, 15A18, 05C31}
\keywords{spectral radius, nonregular graph, prescribed maximum degree, degree
sequence, extremal graph, Sylvester's criterion, counterexample}

\hypersetup{
  pdftitle={The Liu-Li predicted degree sequence is not extremal at (n,Delta)=(8,5)}
}

\begin{document}

\begin{abstract}
Let $\mathcal{G}(n,\Delta)$ be the set of graphs attaining the maximum spectral
radius among all connected nonregular graphs on $n$ vertices with maximum degree
$\Delta$. Liu and Li conjectured in 2008 that for $3\le\Delta\le n-2$ every
member of $\mathcal{G}(n,\Delta)$ has degree sequence
$(\Delta,\dots,\Delta,\delta)$ with $\delta=\Delta-1$ when $n\Delta$ is odd and
$\delta=\Delta-2$ when $n\Delta$ is even. We exhibit an $18$-edge graph $H$ on
$8$ vertices with $\lambda_1(H)>49/10$ and prove, by an exact integer
certificate, that every connected graph on $8$ vertices with the predicted
degree sequence $(5,5,5,5,5,5,5,3)$ satisfies $\lambda_1<97/20$; that class has
exactly three isomorphism types, and the accompanying program also checks all
$21{,}000$ of its labeled members one by one. Hence no member of
$\mathcal{G}(8,5)$ carries the predicted degree sequence and the conjecture is
false at $(n,\Delta)=(8,5)$. We do \emph{not} claim priority for refuting the
conjecture: counterexamples at $\Delta=n-2$ are an undrawn corollary of a 2024
preprint of Huang, Liu and Yang, as we explain in Section~\ref{sec:priority}.
\end{abstract}

\maketitle

\section{Prior work}\label{sec:priority}

The bare assertion ``the Liu--Li 2008 conjecture is false'' is not new. Huang,
Liu and Yang \cite{HLY} determine $\mathcal{G}(n,n-2)$ exactly (their
Theorem~7), which yields counterexamples to the conjecture at every even
$n\ge 8$ with $\Delta=n-2$, and determine $\mathcal{G}(n,n-3)$ for $n\ge 59$
(their Theorem~8), which yields counterexamples at every even $n\ge 60$ with
$\Delta=n-3$. They state the Liu--Li conjecture as their own Conjecture~1 and do
not remark that their theorems contradict it. We claim no priority for the
refutation.

What is offered here is one cell, $(n,\Delta)=(8,5)$. It has $\Delta=n-3$ and
$n=8$, so it lies in the gap between the two theorems of \cite{HLY}, and Liu
\cite{Liu2024} proves the conjecture for $\Delta=3$ and $\Delta=4$; the argument
below is an exact, solver-free integer certificate over the \emph{complete}
predicted class. Whether the cell was already settled by
Table~1 of \cite{LSW2007}, which we could not read, is discussed in
Section~\ref{sec:notsettled}.

\section{The conjecture, verbatim}\label{sec:conj}

Let $\lambda_1(G)$ denote the largest eigenvalue of the adjacency matrix of $G$.
Following \cite[Section on extreme eigenvalues of nonregular
graphs]{LiuNing}: ``Let $\mathcal{G}(n,\Delta)$ denote the set of graphs
attaining the maximum spectral radius among all connected nonregular graphs with
$n$ vertices and maximum degree $\Delta$, and let $\lambda(n,\Delta)$ denote the
maximum spectral radius.''

The target is the following. It appears as Conjecture~2.12 of Liu and Li
\cite{LiuLi}, in the form ``Let $G$ be a connected non-regular graph on $n$
vertices and $2<\Delta<n-1$. Then $G$ is $\lambda_1$-extremal if and only if $G$
is a graph with the greatest maximum eigenvalue in classes $C_\pi$ and
$\pi=(\Delta,\Delta,\dots,\Delta,\delta)$, where $\delta=(\Delta-1)$, when
$n\Delta$ is odd, $\delta=(\Delta-2)$, when $n\Delta$ is even.'' Two modern
restatements keep only the forward direction, which is the weaker statement:
verbatim from the e-print of \cite{LiuNing} and, character for character
identically, from the e-print of \cite{Liu2024}:

\begin{quote}
Let $3\leq\Delta\leq n-2$ and $G\in\mathcal{G}(n,\Delta)$. Then $G$ has degree
sequence $(\Delta,\ldots,\Delta,\delta)$, where
\[
\delta =
\begin{cases}
\Delta-1, & n\Delta\ \text{is odd},\\
\Delta-2, & n\Delta\ \text{is even}.
\end{cases}
\]
\end{quote}

The quantification, made explicit with no words added to the mathematics: for
all integers $n,\Delta$ with $3\le\Delta\le n-2$ and for \emph{all}
$G\in\mathcal{G}(n,\Delta)$ the degree sequence of $G$ is
$(\Delta,\dots,\Delta,\delta)$ with $\delta$ as displayed. A counterexample to
this forward-only form is a counterexample to the ``if and only if'' of
\cite{LiuLi}, which is the stronger statement.

\section{The cell \texorpdfstring{$(n,\Delta)=(8,5)$}{(n,Delta)=(8,5)}}
\label{sec:cell}

The ambient clause holds: $3\le 5\le n-2=6$. Since $n\Delta=40$ is even, the
rule predicts $\delta=\Delta-2=3$, i.e.\ degree sequence
$(5,5,5,5,5,5,5,3)$, written $(5^7,3)$ below. Such a graph has degree sum
$7\cdot 5+3=38$, hence $19$ edges, and its complement on $8$ vertices has
$28-19=9$ edges with degree sequence $(4,2,2,2,2,2,2,2)$. Since $7\cdot5=35$ is
odd, any $\delta$ with $(5^7,\delta)$ graphical must be odd, and nonregularity
forces $1\le\delta\le 4$; so the only available sequences of the shape ``all but
one vertex of degree $\Delta$'' are $(5^7,3)$ and $(5^7,1)$.

Liu, Huang and You \cite{LHY2009} write in their Remark~2.11:
``Conjecture~1.4 is true for small $n$, where $n\le 7$. Since $2<\Delta<n-1$, it
need to be verified for $n=5,6,7$ as follows'', and then list exactly the six
cells $g(5,3),g(6,3),g(6,4),g(7,3),g(7,4),g(7,5)$ with their extremal degree
sequences and spectral radii; that verification stops at $n=7$.

\section{The exhibited objects}

\subsection{The witness \texorpdfstring{$H$}{H}}
On the vertex set $\{0,\dots,7\}$ let $H$ have the $18$ edges
\[
\begin{array}{l}
07,\;12,\;13,\;14,\;16,\;17,\;24,\;25,\;26,\;27,\\
34,\;35,\;36,\;37,\;45,\;46,\;56,\;57 .
\end{array}
\]
Then $\deg(0)=1$ and $\deg(v)=5$ for $v=1,\dots,7$, so the degree sum is
$36=2\cdot 18$; the degree sequence is $(5^7,1)$; the maximum degree is exactly
$5$; $H$ is not regular; and $H$ is connected, since $\{1,\dots,7\}$ induces
$K_7$ minus four edges and vertex $0$ hangs off vertex $7$. In graph6 this
labelling is \texttt{GI[z\textasciitilde{}g}.

The four pairs missing from the core $\{1,\dots,7\}$ are $15,23,47,67$, i.e.\ the
path $4$--$7$--$6$ together with two disjoint edges: $P_3+2K_2$. The pendant
vertex $0$ is attached to vertex $7$, the centre of that $P_3$. So
\[
H=\bigl(K_7-(P_3+2K_2)\bigr)\ \text{with a pendant vertex on the $P_3$ centre}.
\]

A second labelling of the same graph, graph6 \texttt{GNz\textasciitilde{}s?}, is
used below because the Rayleigh vector is written in it. Its $18$ edges are
\[
\begin{array}{l}
03,\;04,\;05,\;06,\;07,\;12,\;13,\;14,\;15,\;16,\\
23,\;24,\;25,\;26,\;35,\;36,\;45,\;46 .
\end{array}
\]
The two labellings are cospectral, and the characteristic polynomial of $H$ is
\begin{equation}\label{eq:pH}
\begin{aligned}
p_H(x)&=x^8-18x^6-32x^5+5x^4+32x^3+12x^2\\
      &=x^2\bigl(x^6-18x^4-32x^3+5x^2+32x+12\bigr).
\end{aligned}
\end{equation}

\begin{remark}\label{rem:notafind}
$H$ is exhibited, but it is not a discovery. There are exactly $2{,}520$ labeled
graphs on $8$ vertices with degree sequence $(5^7,1)$, all connected, and they
form a \emph{single} isomorphism class with $|\Aut|=16$, since
$2520=8!/16$. Indeed the complement of such a graph has $10$ edges and degrees
$(6,2^7)$, so it consists of three cycles glued at the single vertex of degree
$6$; on $8$ vertices the only possible lengths are $3,3,4$, and the labeled
count is $8\cdot\binom{7}{3}\cdot\frac{3!}{2}\cdot 3=2520$. Hence $H$ is
\emph{forced} by its degree sequence, and it is the $n=8$ member $H_1(8)$ of the
family in \cite{HLY}. What is proved below is not a property of the graph but a
comparison: this forced graph has strictly larger spectral radius than every
graph of the degree sequence the conjecture predicts at $(8,5)$. Optimality of
$H$ in $\mathcal{G}(8,5)$ is not proved here.
\end{remark}

\subsection{The predicted class}
The connected graphs on $8$ vertices with degree sequence $(5^7,3)$ number
exactly $21{,}000$ labeled graphs and fall into exactly three isomorphism
classes, with automorphism group orders $48,4,4$ and orbit sizes
$840+10080+10080=21000$. The classification is by hand: the complement has $9$
edges and degrees $(4,2^7)$, so the component containing the degree-$4$ vertex
has $k+1$ edges on $k$ vertices with all degrees in $\{2,4\}$ and a single $4$,
hence is two cycles glued at one vertex (a theta graph would need two vertices
of degree $3$); the remaining vertices form disjoint cycles. On $8$ vertices with
$9$ edges the only possibilities for the glued lengths, together with the
leftover cycles, are $(3,3)$ plus a disjoint $C_3$; $(3,6)$; and $(4,5)$, since
a leftover of one or two vertices cannot form a cycle. The corresponding labeled
counts are $8\binom{7}{4}\cdot 3=840$, $8\binom{7}{2}\cdot\frac{5!}{2}=10080$ and
$8\binom{7}{3}\cdot\frac{3!}{2}\cdot\frac{4!}{2}=10080$, which sum to exactly the
labeled total, so no class is missed and none is double counted.

Representatives, each with $19$ edges:
\[
\begin{array}{l}
C_1\ (\texttt{GFz\textasciitilde{}v?}):\ \,
03,04,05,06,07,13,14,15,16,17,\\
\qquad\qquad\qquad\ \ 23,24,25,26,27,35,36,45,46\\[3pt]
C_2\ (\texttt{GLv\textasciitilde{}v?}):\ \,
03,04,05,06,07,12,14,15,16,17,\\
\qquad\qquad\qquad\ \ 23,25,26,27,34,35,36,45,46\\[3pt]
C_3\ (\texttt{GZn]\textasciitilde{}?}):\ \,
02,04,05,06,07,12,13,15,16,17,\\
\qquad\qquad\qquad\ \ 23,24,27,34,35,36,45,46,56
\end{array}
\]
with characteristic polynomials, as coefficient lists from $x^8$ down to $x^0$,
\begin{equation}\label{eq:pC}
\begin{aligned}
p_{C_1}&:\ [1,0,-19,-24,12,0,0,0,0],\\
p_{C_2}&:\ [1,0,-19,-28,24,36,-8,0,0],\\
p_{C_3}&:\ [1,0,-19,-30,28,62,-3,-24,1].
\end{aligned}
\end{equation}
Which representative realises which complement type is recorded in the
accompanying program and is not needed below.

\section{The refutation}

The argument is three exact facts and one finiteness remark. Nothing
floating point enters it.

\begin{lemma}[membership]\label{lem:member}
$H$ is connected, nonregular, has maximum degree exactly $5$, and lies on $n=8$
vertices with $3\le 5\le 6$. Hence $H$ belongs to the family over which
$\mathcal{G}(8,5)$ is defined.
\end{lemma}

\begin{proof}
Immediate from the edge list and the degrees in Section~4.1.
\end{proof}

\begin{lemma}[lower bound]\label{lem:lower}
$\lambda_1(H)>49/10$.
\end{lemma}

\begin{proof}
In the second labelling of Section~4.1, with $A$ the adjacency matrix, the explicit
integer vector
\[
\begin{aligned}
x=(\,&3251590,\,3914936,\,3914936,\,3818824,\\
    &3818824,\,3818824,\,3818824,\,663346\,)
\end{aligned}
\]
has exact rational Rayleigh quotient
\[
\frac{x^{\mathsf T}Ax}{x^{\mathsf T}x}
=\frac{61272483059723}{12499997522989}=4.901799616\ldots>\frac{49}{10}.
\]
Since $A$ is real symmetric, $\lambda_1(A)=\max_{x\ne 0}x^{\mathsf T}Ax/x^{\mathsf T}x$,
so $\lambda_1(H)>49/10$. Independently, $p_H$ of \eqref{eq:pH} is monic with all
roots real and $p_H(49/10)=-27065510199/100000000<0$, which forces its largest
root above $49/10$.
\end{proof}

\begin{lemma}[upper bound over the whole predicted class]\label{lem:upper}
Every connected graph $G$ on $8$ vertices with degree sequence $(5^7,3)$
satisfies $\lambda_1(G)<97/20$.
\end{lemma}

\begin{proof}
For a symmetric matrix $A$ and integers $c,d>0$, the matrix $cI-dA$ is positive
definite if and only if $\lambda_1(A)<c/d$; and by Sylvester's criterion
positive definiteness is equivalent to the positivity of all $n$ leading
principal minors, which for an integer matrix is a statement about $n$ integers.

For the three representatives the eight leading principal minors of the integer
matrix $97I-20A$ are
\[
\begin{array}{ll}
C_1:&[\,97,\ 9409,\ 912673,\ 77238481,\ 6396925057,\ 355826560529,\\
    &\ \ 8841684881313,\ 28053247092561\,]\\[2pt]
C_2:&[\,97,\ 9409,\ 873873,\ 77398481,\ 6308461057,\ 349728752529,\\
    &\ \ 8491067905313,\ 23557099988561\,]\\[2pt]
C_3:&[\,97,\ 9409,\ 835073,\ 72082881,\ 5633183857,\ 364595360529,\\
    &\ \ 8202660545313,\ 18809429316561\,]
\end{array}
\]
all positive, so $\lambda_1(C_i)<97/20$ for $i=1,2,3$. The last entry of each
list cross-checks against \eqref{eq:pC}, since
$\det(97I-20A)=20^8\,p(97/20)$ with $20^8=25{,}600{,}000{,}000$. Since spectra
are isomorphism invariants and the class has
exactly the three isomorphism classes classified in Section~4.2, the bound holds
for every member.
\end{proof}

\begin{remark}[the classification is not on the critical path]\label{rem:sweep}
Lemma~\ref{lem:upper} as proved above rests on the hand classification of
Section~4.2, i.e.\ on a completeness argument. The accompanying program
\texttt{verify.py} removes that dependency: it generates all $21{,}000$ labeled
graphs of degree sequence $(5^7,3)$ directly, by enumerating their $9$-edge
complements of degrees $(4,2^7)$, confirms that all $21{,}000$ are connected, and
verifies Sylvester's criterion on $97I-20A$ for each one individually, with zero
violations.
\end{remark}

\begin{theorem}\label{thm:main}
The Liu--Li conjecture is false at $(n,\Delta)=(8,5)$: no member of
$\mathcal{G}(8,5)$ has degree sequence $(5,5,5,5,5,5,5,3)$.
\end{theorem}

\begin{proof}
The family of connected nonregular graphs on $8$ vertices with maximum degree
exactly $5$ is finite and, by Lemma~\ref{lem:member}, non-empty, so
$\mathcal{G}(8,5)$ is non-empty. Let $M\in\mathcal{G}(8,5)$. By maximality and
Lemmas~\ref{lem:member}, \ref{lem:lower} and \ref{lem:upper},
\[
\lambda_1(M)\ \ge\ \lambda_1(H)\ >\ \frac{49}{10}\ >\ \frac{97}{20}\ >\ \lambda_1(G)
\]
for every connected $G$ on $8$ vertices with degree sequence $(5^7,3)$. Hence
$M$ does not have that degree sequence. Since $n\Delta=40$ is even, the
conjecture demands $\delta=\Delta-2=3$, i.e.\ exactly that sequence, for every
member of $\mathcal{G}(8,5)$. The conjecture therefore fails at $(8,5)$, and
since the pinned form is the forward direction of the ``if and only if'' of
\cite{LiuLi}, the original fails there too.
\end{proof}

\begin{remark}[the sequence carried by the witness]
$H$ has degree sequence $(5^7,1)$, i.e.\ $\delta=1=\Delta-4$. The replacement
conjecture of Liu and Ning \cite{LiuNing} (Conjecture~2 of \cite{HLY}) predicts
$\delta=1$ for $\Delta$ odd and $n$ even, and $\delta=1$ is also the shape that
Theorem~8 of \cite{HLY} proves asymptotically at $\Delta=n-3$; the witness here
is consistent with that prediction and does not prove it.
\end{remark}

\section{Scope}\label{sec:notsettled}

Only the cell $(8,5)$ is reported. The refutation of the conjecture is not new
and no priority is claimed here (Section~\ref{sec:priority}); the increment is
this one cell together with an exact whole-class certificate. The witness is not
a find: it is forced by its degree sequence and it \emph{is} $H_1(8)$ of
\cite{HLY} (Remark~\ref{rem:notafind}). Which graph attains $\lambda(8,5)$ is
not established here, and Theorem~\ref{thm:main} does not need it: it needs only
that $H$ beats every graph of the predicted degree sequence. In particular we do
not assert as a certified statement that $H\in\mathcal{G}(8,5)$, and Liu's
replacement conjecture is not proved either. Nothing here concerns any other $n$
or $\Delta$: the $\Delta=n-2$ cells, including at $n=8$, are Theorem~7 of
\cite{HLY}, and at $\Delta=5$ the cells $n=10$ and $n=12$ are untouched.

One residual priority risk remains. We could not obtain Table~1 of
\cite{LSW2007}, the table of $\lambda_1$-extremal graphs that \cite{LiuLi}
explicitly consults (``by checking the Table 1 of [7]''); if it reaches $n=8$,
this cell was settled in 2007. The indirect evidence against is that Bolian Liu,
an author of both \cite{LiuLi} and \cite{LHY2009}, wrote in 2009 that the
conjecture ``need[s] to be verified for $n=5,6,7$'' and listed exactly those six
cells, and that both \cite{Liu2024} and \cite{HLY} treat $\Delta\ge 5$ as open.
A successor should obtain that table.

\section*{Verification}

The accompanying program \texttt{verify.py} (Python 3.9+, standard library only)
re-derives the quantities transcribed into it from the objects printed above: the
degrees, connectivity, structure and graph6 encodings of $H$ and of
$C_1,C_2,C_3$; the characteristic polynomials, by Newton's identities on
integer traces; the Rayleigh quotient; the exact rational evaluations; and the
leading principal minors of $97I-20A$. It does not parse this paper, so it does
not certify that that list of transcribed quantities is complete. It then carries
out the complete enumerations of Remarks~\ref{rem:notafind} and \ref{rem:sweep}
and runs certifier controls in both polarities. It does not re-derive the
statement of the Liu--Li conjecture, which it takes from Section~\ref{sec:conj},
and it does not establish which graph attains $\lambda(8,5)$. Its recorded output
accompanies this paper.

\begin{thebibliography}{9}

\bibitem{HLY}
Z.~Huang, J.~Liu and C.~Yang,
\emph{Nonregular graphs with a given maximum degree attaining maximum spectral
radius},
\href{https://arxiv.org/abs/2411.17371}{arXiv:2411.17371}, 26 November 2024.

\bibitem{LiuLi}
B.~Liu and G.~Li,
\emph{A note on the largest eigenvalue of non-regular graphs},
Electron. J. Linear Algebra \textbf{17} (2008), 54--61.
\href{https://doi.org/10.13001/1081-3810.1249}{doi:10.13001/1081-3810.1249}

\bibitem{LHY2009}
B.~Liu, Y.~Huang and Z.~You,
\emph{On $\lambda_1$-extremal non-regular graphs},
Electron. J. Linear Algebra \textbf{18} (2009), 735--744.
\href{https://doi.org/10.13001/1081-3810.1341}{doi:10.13001/1081-3810.1341}

\bibitem{Liu2024}
L.~Liu,
\emph{Extremal spectral radius of nonregular graphs with prescribed maximum
degree},
J. Combin. Theory Ser. B \textbf{169} (2024), 430--479;
\href{https://arxiv.org/abs/2203.10245}{arXiv:2203.10245}.

\bibitem{LiuNing}
L.~Liu and B.~Ning,
\emph{Unsolved problems in spectral graph theory},
Oper. Res. Trans. \textbf{27} (2023), no.~4, 33--60;
\href{https://arxiv.org/abs/2305.10290}{arXiv:2305.10290}.

\bibitem{LSW2007}
B.~Liu, J.~Shen and X.~Wang,
\emph{On the largest eigenvalue of non-regular graphs},
J. Combin. Theory Ser. B \textbf{97} (2007), 1010--1018.
\href{https://doi.org/10.1016/j.jctb.2007.02.008}{doi:10.1016/j.jctb.2007.02.008}

\end{thebibliography}

\end{document}
